\subsection{Topic 4: Compilation, Interpretation and Subroutines – Exercise Sheet [Solutions]}

\subsection*{Q\#1: A letter means push and an asterisk means pop. Show the output of the pop operations:}

\begin{itemize}
\item[a.] M M I B * * R * * G N I * * * A H * * * \\
\textbf{Output:} BIRMINGHAM

\item[b.] A U E * * I R * * K * * \\
\textbf{Output:} EURIKA
\end{itemize}

\subsection*{Q\#2: Convert infix expressions to postfix (RPN) using Shunting Yard algorithm.}
\textbf{a. } $5 - 15 \times (25 + 5)$

\begin{tabularx}{\textwidth}{|c|X|X|X|}
\hline
\textbf{Step \#} & \textbf{Output} & \textbf{Stack (Top on Right)} & \textbf{Input} \\
\hline
1 & Empty & Empty & $5 - 15 \times (25 + 5)$ \\
\hline
2 & 5 & Empty & $- 15 \times (25 + 5)$ \\
\hline
3 & 5 & – & $15 \times (25 + 5)$ \\
\hline
4 & 5 15 & – & $\times (25 + 5)$ \\
\hline
5 & 5 15 & – * & $(25 + 5)$ \\
\hline
6 & 5 15 & – * ( & $25 + 5)$ \\
\hline
7 & 5 15 25 & – * ( & $+ 5)$ \\
\hline
8 & 5 15 25 & – * ( + & 5) \\
\hline
9 & 5 15 25 5 & – * ( + & ) \\
\hline
10 & 5 15 25 5 + * – & Empty & Empty \\
\hline
\end{tabularx}

\vspace{1em}

\textbf{b. } $4 \times (6 + 10 / 2.5 - 8)$

\begin{tabularx}{\textwidth}{|c|X|X|X|}
\hline
\textbf{Step \#} & \textbf{Output} & \textbf{Stack (Top on Right)} & \textbf{Input} \\
\hline
1 & Empty & Empty & $4 \times (6 + 10 / 2.5 - 8)$ \\
\hline
2 & 4 & Empty & $\times (6 + 10 / 2.5 - 8)$ \\
\hline
3 & 4 & \(\times\) & $(6 + 10 / 2.5 - 8)$ \\
\hline
4 & 4 & \(\times\) ( & $6 + 10 / 2.5 - 8)$ \\
\hline
5 & 4 6 & \(\times\) ( & $+ 10 / 2.5 - 8)$ \\
\hline
6 & 4 6 & \(\times\) ( + & $10/ 2.5 - 8)$ \\
\hline
7 & 4 6 10 & \(\times\) ( + & $/2.5 - 8)$ \\
\hline
8 & 4 6 10 & \(\times\) ( + / & $2.5- 8)$ \\
\hline
9 & 4 6 10 2.5 & \(\times\) ( + / & $- 8)$ \\
\hline
10 & 4 6 10 2.5 / + & \(\times\) ( – & $8)$ \\
\hline
11 & 4 6 10 2.5 / + 8 & \(\times\) ( – & ) \\
\hline
12 & 4 6 10 2.5 / + 8 – & \(\times\) & Empty \\
\hline
13 & 4 6 10 2.5 / + 8 – \(\times\) & Empty & Empty \\
\hline
\end{tabularx}

\vspace{1em}

\textbf{c. } $17 + 21 / 3 \times 7 + 21$

\begin{center}
\begin{tabular}{|c|l|l|l|}
\hline
\textbf{Step \#} & \textbf{Output} & \textbf{Stack (Top on Right)} & \textbf{Input} \\
\hline
1 & Empty & Empty & $17 + 21 / 3 \times 7 + 21$ \\
2 & 17 & Empty & $+ 21 / 3 \times 7 + 21$ \\
3 & 17 & + & $21 / 3 \times 7 + 21$ \\
4 & 17 21 & + & $/ 3 \times 7 + 21$ \\
5 & 17 21 3 & + / & $\times 7 + 21$ \\
6 & 17 21 3 / & + & $\times 7 + 21$ \\
7 & 17 21 3 / 7 & + * & $+ 21$ \\
8 & 17 21 3 / 7 * & + & $21$ \\
9 & 17 21 3 / 7 * + 21 & + & Empty \\
10 & 17 21 3 / 7 * + 21 + & Empty & Empty \\
\hline
\end{tabular}
\end{center}

\subsection*{Q\#3: Evaluate the RPN expressions obtained in the question above, using a stack to check that they
calculate the correct answers.}
\textbf{a. RPN: } $5\ 15\ 25\ 5\ +\ *\ -$

\begin{center}
\begin{tabularx}{\textwidth}{|c|X|X|}
\hline
\textbf{Step \#} & \textbf{Input} & \textbf{Stack (Top on Right)} \\
\hline
1 & 5 & 5 \\
\hline
2 & 15 & 5 15 \\
\hline
3 & 25 & 5 15 25 \\
\hline
4 & 5 & 5 15 25 5 \\
\hline
5 & + & 5 15 30 \\
\hline
6 & * & 5 450 \\
\hline
7 & – & \textbf{-445} \\
\hline
\end{tabularx}
\end{center}

\vspace{1em}

\textbf{b. RPN: } $4\ 6\ 10\ 2.5\ /\ +\ 8\ -\ *$

\begin{center}
\begin{tabular}{|c|c|l|}
\hline
\textbf{Step \#} & \textbf{Input} & \textbf{Stack (Top on Right)} \\
\hline
1 & 4 & 4 \\
2 & 6 & 4 6 \\
3 & 10 & 4 6 10 \\
4 & 2.5 & 4 6 10 2.5 \\
5 & / & 4 6 4 \\
6 & + & 4 10 \\
7 & 8 & 4 10 8 \\
8 & – & 4 2 \\
9 & * & \textbf{8} \\
\hline
\end{tabular}
\end{center}

\vspace{1em}

\textbf{c. RPN: } $17\ 21\ 3\ /\ 7\ *\ +\ 21\ +$

\begin{center}
\begin{tabular}{|c|c|l|}
\hline
\textbf{Step \#} & \textbf{Input} & \textbf{Stack (Top on Right)} \\
\hline
1 & 17 & 17 \\
2 & 21 & 17 21 \\
3 & 3 & 17 21 3 \\
4 & / & 17 7 \\
5 & 7 & 17 7 7 \\
6 & * & 17 49 \\
7 & + & 66 \\
8 & 21 & 66 21 \\
9 & + & \textbf{87} \\
\hline
\end{tabular}
\end{center}

\subsection*{Q\#4:}
Convert Java: \texttt{i = (j + k) \&\& m}

\begin{center}
\begin{tabular}{|c|c|c|c|}
\hline
\textbf{Java} & \textbf{IJVM Bytecode} & \textbf{Hex Bytecode} & \textbf{Stack (Top on Right)} \\
\hline
ILOAD j & ILOAD & 0x15 0x02 & [17] \\
ILOAD k & ILOAD & 0x15 0x03 & [17, 14] \\
IADD & IADD & 0x60 & [31] \\
ILOAD m & ILOAD & 0x15 0x04 & [31, 45] \\
IAND & IAND & 0x7E & [1] \\
ISTORE i & ISTORE & 0x36 0x01 & [] \\
\hline
\end{tabular}
\end{center}

\subsection*{Q\#5:}
Bytecode for $x = (p \times q) - (r + s)$, given p=5, q=4, r=2, s=1:

\begin{center}
\begin{tabular}{|c|l|}
\hline
\textbf{IJVM Bytecode} & \textbf{Stack (Top on Right)} \\
\hline
ILOAD p & [5] \\
ILOAD q & [5, 4] \\
IMUL & [20] \\
ILOAD r & [20, 2] \\
ILOAD s & [20, 2, 1] \\
IADD & [20, 3] \\
ISUB & [17] \\
ISTORE x & [] \\
\hline
\end{tabular}
\end{center}

\subsection*{Q\#6 (Optional):} A palindrome is a word, phrase, number, or sequence of characters that reads the
same backward as forward. In other words, if you reverse the order of the characters, the resulting
string is identical to the original. Palindromes are usually case-insensitive (e.g., "Racecar" is considered
a palindrome). Examples include "level", "radar", "civic", 121, 12321.
To detect if a given string is a palindrome using a stack, you can utilize the stack's Last In First Out
(LIFO) property. The basic idea is to push the characters of the string onto a stack and then pop them
off to compare with the original string. Write a program in Python to detect if the given input string is a
palindrome.

\begin{lstlisting}[language=Python]
def is_palindrome(input_string):
    stack = []

    for char in input_string:
        stack.append(char)

    for char in input_string:
        if stack.pop() != char:
            return False
    return True

# Test the function
test_string = "radar"
if is_palindrome(test_string):
    print(f'"{test_string}" is a palindrome.')
else:
    print(f'"{test_string}" is not a palindrome.')
\end{lstlisting}

% \end{document}
